% GENERATED FILE. Edit canonical lessons, then run npm run generate:books.
% canonical-source-hash: fnv1a64:c4ddcca2b92ed9a9
% canonical-lessons: GU-C543-apvad, GU-C543-yogya, GU-C543-layak, GU-C543-kabel, GU-C543-khas

\chapter{Fit for the Job}
\label{ch:gu-fit-for-the-job}
\hlchaptermodality{543}

\begin{chapteropening}
\textbf{By the end of this chapter:} \emph{I can say who is fit for something and when an exception is made.}

Complete the last lesson of chapter 543: I can say who is fit for something and when an exception is made.
\end{chapteropening}

\section[\texorpdfstring{\gu{અપવાદ} (apvād)}{અપવાદ (apvād)}]{\gu{અપવાદ} (apvād) — an exception}

\label{lesson:GU-C543-apvad}



\begin{quote}
Before the new one: say the Gujarati for a recommendation, then the Gujarati for a polite request.
\end{quote}

\subsection*{You'll want to know: \gu{અપવાદ}}
\textbf{\gu{અપવાદ}} — \emph{apvād} — \textquotedblleft{}an exception\textquotedblright{}.

\begin{quote}
\gu{જો} \gu{તમે} \gu{બીમાર} \gu{હશો} \gu{તો} \gu{અમે} \gu{અપવાદ} \gu{કરીશું}.
\end{quote}

\begin{quote}
\emph{jo tame bīmār haśo to ame apvād karīśũ.}
\end{quote}

\begin{quote}
If you are ill, we will make an exception.
\end{quote}

A masculine noun.

\begin{cousinweb}
\textbf{\gu{અપવાદ}} (\emph{apvād}) is a tatsama, taken from Sanskrit \emph{apavāda}, \textquotedblleft{}speaking against; a refutation\textquotedblright{} (Monier-Williams); in grammar and in rules it is the case that does not follow.
\end{cousinweb}

\begin{grammarlens}[title={What you've built}]
One more word for saying what happens if.
\end{grammarlens}

\subsection*{Your turn}
\begin{itemize}
\raggedright
  \item \emph{Say it:} \emph{apvād}
  \item \emph{Say it:} \emph{apvād}, once more
  \item \emph{Say it:} \emph{jo tame bīmār haśo to ame apvād karīśũ.} — If you are ill, we will make an exception.
  \item \emph{Recall:} say the Gujarati for a recommendation, then the Gujarati for a polite request, then say \emph{apvād} again
\end{itemize}

\begin{tcolorbox}[breakable,colback=teal!4,colframe=teal!35!black,title={Before you move on}]
What does \gu{અપવાદ} mean? (\textquotedblleft{}An exception\textquotedblright{}.) Say it once more.
\end{tcolorbox}
\section[\texorpdfstring{\gu{યોગ્ય} (yogya)}{યોગ્ય (yogya)}]{\gu{યોગ્ય} (yogya) — suitable}

\label{lesson:GU-C543-yogya}



\begin{quote}
Before the new one: say the Gujarati for a polite request, then the Gujarati for an exception.
\end{quote}

\subsection*{You'll want to know: \gu{યોગ્ય}}
\textbf{\gu{યોગ્ય}} — \emph{yogya} — \textquotedblleft{}suitable\textquotedblright{}.

\begin{quote}
\gu{જો} \gu{આ} \gu{છોકરો} \gu{યોગ્ય} \gu{હશે} \gu{તો} \gu{અમે} \gu{તેને} \gu{નોકરી} \gu{આપીશું}.
\end{quote}

\begin{quote}
\emph{jo ā chokro yogya haśe to ame tene nokrī āpīśũ.}
\end{quote}

\begin{quote}
If this boy is suitable, we will give him the job.
\end{quote}

\begin{cousinweb}
\textbf{\gu{યોગ્ય}} (\emph{yogya}) is a tatsama, taken from Sanskrit \emph{yogya}, \textquotedblleft{}fit for the yoke; fit, suitable\textquotedblright{} (Monier-Williams).
\end{cousinweb}

\begin{grammarlens}[title={What you've built}]
One more word for saying what happens if.
\end{grammarlens}

\subsection*{Your turn}
\begin{itemize}
\raggedright
  \item \emph{Say it:} \emph{yogya}
  \item \emph{Say it:} \emph{yogya}, once more
  \item \emph{Say it:} \emph{jo ā chokro yogya haśe to ame tene nokrī āpīśũ.} — If this boy is suitable, we will give him the job.
  \item \emph{Recall:} say the Gujarati for a polite request, then the Gujarati for an exception, then say \emph{yogya} again
\end{itemize}

\begin{tcolorbox}[breakable,colback=teal!4,colframe=teal!35!black,title={Before you move on}]
What does \gu{યોગ્ય} mean? (\textquotedblleft{}Suitable\textquotedblright{}.) Say it once more.
\end{tcolorbox}
\section[\texorpdfstring{\gu{લાયક} (lāyak)}{લાયક (lāyak)}]{\gu{લાયક} (lāyak) — deserving}

\label{lesson:GU-C543-layak}



\begin{quote}
Before the new one: say the Gujarati for an exception, then the Gujarati for suitable.
\end{quote}

\subsection*{You'll want to know: \gu{લાયક}}
\textbf{\gu{લાયક}} — \emph{lāyak} — \textquotedblleft{}deserving\textquotedblright{}.

\begin{quote}
\gu{જો} \gu{તમે} \gu{લાયક} \gu{હશો} \gu{તો} \gu{તમને} \gu{બઢતી} \gu{મળશે}.
\end{quote}

\begin{quote}
\emph{jo tame lāyak haśo to tamne baḍhtī maḷśe.}
\end{quote}

\begin{quote}
If you deserve it, you will get a promotion.
\end{quote}

\begin{cousinweb}
\textbf{\gu{લાયક}} (\emph{lāyak}) is usually traced to Arabic \emph{lāyiq}, \textquotedblleft{}worthy, fitting\textquotedblright{}, through Persian. The quality it names, \textbf{\gu{લાયકાત}}, is a qualification.
\end{cousinweb}

\begin{grammarlens}[title={What you've built}]
One more word for saying what happens if.
\end{grammarlens}

\subsection*{Your turn}
\begin{itemize}
\raggedright
  \item \emph{Say it:} \emph{lāyak}
  \item \emph{Say it:} \emph{lāyak}, once more
  \item \emph{Say it:} \emph{jo tame lāyak haśo to tamne baḍhtī maḷśe.} — If you deserve it, you will get a promotion.
  \item \emph{Recall:} say the Gujarati for an exception, then the Gujarati for suitable, then say \emph{lāyak} again
\end{itemize}

\begin{tcolorbox}[breakable,colback=teal!4,colframe=teal!35!black,title={Before you move on}]
What does \gu{લાયક} mean? (\textquotedblleft{}Deserving\textquotedblright{}.) Say it once more.
\end{tcolorbox}
\section[\texorpdfstring{\gu{કાબેલ} (kābel)}{કાબેલ (kābel)}]{\gu{કાબેલ} (kābel) — capable}

\label{lesson:GU-C543-kabel}



\begin{quote}
Before the new one: say the Gujarati for suitable, then the Gujarati for deserving.
\end{quote}

\subsection*{You'll want to know: \gu{કાબેલ}}
\textbf{\gu{કાબેલ}} — \emph{kābel} — \textquotedblleft{}capable\textquotedblright{}.

\begin{quote}
\gu{જો} \gu{મિસ્ત્રી} \gu{કાબેલ} \gu{હશે} \gu{તો} \gu{ગાડી} \gu{એક} \gu{કલાકમાં} \gu{ઠીક} \gu{થઈ} \gu{જશે}.
\end{quote}

\begin{quote}
\emph{jo mistrī kābel haśe to gāḍī ek kalākmā̃ ṭhīk thaī jaśe.}
\end{quote}

\begin{quote}
If the mechanic is capable, the car will be fixed in an hour.
\end{quote}

\begin{cousinweb}
\textbf{\gu{કાબેલ}} (\emph{kābel}) is usually traced to Arabic \emph{qābil}, \textquotedblleft{}able, capable\textquotedblright{}, through Persian.
\end{cousinweb}

\begin{grammarlens}[title={What you've built}]
One more word for saying what happens if.
\end{grammarlens}

\subsection*{Your turn}
\begin{itemize}
\raggedright
  \item \emph{Say it:} \emph{kābel}
  \item \emph{Say it:} \emph{kābel}, once more
  \item \emph{Say it:} \emph{jo mistrī kābel haśe to gāḍī ek kalākmā̃ ṭhīk thaī jaśe.} — If the mechanic is capable, the car will be fixed in an hour.
  \item \emph{Recall:} say the Gujarati for suitable, then the Gujarati for deserving, then say \emph{kābel} again
\end{itemize}

\begin{tcolorbox}[breakable,colback=teal!4,colframe=teal!35!black,title={Before you move on}]
What does \gu{કાબેલ} mean? (\textquotedblleft{}Capable\textquotedblright{}.) Say it once more.
\end{tcolorbox}
\section[\texorpdfstring{\gu{ખાસ} (khās)}{ખાસ (khās)}]{\gu{ખાસ} (khās) — special}

\label{lesson:GU-C543-khas}



\begin{quote}
Before the new one: say the Gujarati for deserving, then the Gujarati for capable.
\end{quote}

\subsection*{You'll want to know: \gu{ખાસ}}
\textbf{\gu{ખાસ}} — \emph{khās} — \textquotedblleft{}special\textquotedblright{}.

\begin{quote}
\gu{જો} \gu{કોઈ} \gu{ખાસ} \gu{કારણ} \gu{હશે} \gu{તો} \gu{જ} \gu{તમને} \gu{રજા} \gu{મળશે}.
\end{quote}

\begin{quote}
\emph{jo koī khās kāraṇ haśe to ja tamne rajā maḷśe.}
\end{quote}

\begin{quote}
You'll get leave only if there is a special reason.
\end{quote}

\begin{cousinweb}
\textbf{\gu{ખાસ}} (\emph{khās}) is usually traced to Arabic \emph{khāṣṣ}, \textquotedblleft{}particular, special\textquotedblright{}, through Persian.
\end{cousinweb}

\begin{grammarlens}[title={What you've built}]
One more word for saying what happens if.
\end{grammarlens}

\subsection*{Your turn}
\begin{itemize}
\raggedright
  \item \emph{Say it:} \emph{khās}
  \item \emph{Say it:} \emph{khās}, once more
  \item \emph{Say it:} \emph{jo koī khās kāraṇ haśe to ja tamne rajā maḷśe.} — You'll get leave only if there is a special reason.
  \item \emph{Recall:} say the Gujarati for deserving, then the Gujarati for capable, then say \emph{khās} again
\end{itemize}

\begin{tcolorbox}[breakable,colback=teal!4,colframe=teal!35!black,title={Before you move on}]
What does \gu{ખાસ} mean? (\textquotedblleft{}Special\textquotedblright{}.) Say it once more.
\end{tcolorbox}
